\documentclass[10pt,a4paper]{amsart}
\usepackage[margin=19mm]{geometry}
\usepackage{amsmath,amssymb,mathtools}
\usepackage[scr=rsfs,cal=euler]{mathalfa}
\usepackage{xfrac}
\usepackage[unicode,hidelinks,bookmarks=false]{hyperref}
\newcommand{\ie}{i.e.\ }
\newcommand{\cf}{cf.\ }
\newcommand{\resp}{respectively\ }
\newcommand{\Subsection}{\subsection}
\providecommand{\nobreakdash}{\nobreak\hbox{-}}
\setlength{\emergencystretch}{2em}
\multlinegap0pt
\usepackage{bm}
\usepackage{bbm}
\usepackage{dsfont}
\usepackage{xspace}
\usepackage{cancel}
\usepackage{mathtools}
\usepackage{amsthm}
\makeatletter
\@ifundefined{lemma}{\newtheorem{lemma}{Lemma}}{}
\@ifundefined{proposition}{\newtheorem{proposition}[lemma]{Proposition}}{}
\makeatother
\def\d{\,\mathrm{d}}
\title[Stability and existence of relativistic plasma--vacuum inter]{Stability and existence of relativistic plasma--vacuum interfaces}
\author{Paolo Secchi, Yuri Trakhinin, Tao Wang}
\date{Source: arXiv:2604.26864v1 (CC BY 4.0)}
\begin{document}
\maketitle

\noindent\emph{Source and licence.} Stability and existence of relativistic plasma--vacuum interfaces. Version arXiv:2604.26864v1 (CC BY 4.0), distributed under the Creative Commons Attribution 4.0 International licence, as recorded in the arXiv licence metadata for this exact version and shown on the abstract page https://arxiv.org/abs/2604.26864v1. Full rights notice: licenses/arxiv-2604.26864-CC-BY-4.0.txt.\par\smallskip

\noindent\textit{Editorial scope.} Counted unit: the Lemma whose source label is \texttt{lem:source} in the article (source0.tex lines 3979--3994, source label \texttt{lem:source}), delivered together with the article's own complete original proof of it (source0.tex lines 3995--4045, one \texttt{proof} environment of 51 lines). No mathematics is written, repaired or re-derived: statement and proof are the article's own text line for line. Required context retained uncounted: pro:smooth (lines 3321--3337); smooth.p1 (lines 3324--3334); source (lines 3403--3407); lem:sum2 (lines 3928--3936); es.sum1a (lines 3939--3948). Every definition, display and label of those lines is retained unchanged. Omitted from this package and not redistributed: 1--3320, 3335--3402, 3408--3927, 3937--3938, 3949--3978, 4046--4801. Nothing inside a retained excerpt is omitted or altered: every retained body is delivered exactly as the article's lines are written, so no omission receipt of any kind accompanies this package. Citations: the retained excerpts carry no citation command at all, so no bibliography entry of the article is transcribed and no reference is redistributed. Reference adaptation: none was needed and none was made; every reference inside the delivered unit resolves inside it, so no unresolved reference or citation is produced. Printed slips kept literal: no printed word, symbol, hypothesis or number of the counted statement or of its proof was altered. Layout and class: the article's own class is \texttt{article} with packages including geometry, amsmath, amsfonts, amssymb, bm, mathrsfs, amsthm, amsxtra, xcolor, graphicx, indentfirst, hyperref, natbib, letltxmacro; the wrapper is the lane's pinned \texttt{amsart} editorial wrapper at 10pt on A4, which restates the theorem-like environments the retained text needs and adds the article's own macro definitions that the retained text needs (\texttt{\textbackslash d}), with extra wrapper packages (bm, bbm, dsfont, xspace, cancel, mathtools). The delivered numbering is the unit's own. Assets: no retained span contains a figure environment, a graphics-inclusion command, a PDF-page inclusion, a table or a TikZ picture; no image file is copied into this package, the package declares no asset dependency and no image file is redistributed with it. Licence: version arXiv:2604.26864v1 of this article is distributed under the Creative Commons Attribution 4.0 International licence, as recorded in the arXiv licence metadata for this exact version and shown on the abstract page https://arxiv.org/abs/2604.26864v1. A full rights notice is delivered at licenses/arxiv-2604.26864-CC-BY-4.0.txt. Characters: every retained excerpt is pure ASCII, so the delivered engine, XeLaTeX with its default Latin Modern text font, reports no missing character.
\begin{proposition} \label{pro:smooth}
	Let $T>0$.
	Then there exists a family of smoothing operators $\{\mathcal{S}_{\theta}\}_{\theta\geq 1}$  from $  \mathscr{F}_*^0(\Omega_{T}) $ to $ \bigcap_{k\geq 0}\mathscr{F}_*^k(\Omega_{T})$, such that
	\begin{subequations}\label{smooth.p1}
		\begin{alignat}{2}
			&  \|\mathcal{S}_{\theta} u\|_{H_*^{k}(\Omega_{T})}\lesssim_m \theta^{(k-j)_+}\|u\|_{H_*^{j}(\Omega_{T})}
			&& \textrm{for  \;} k,j\geq 0,    \label{smooth.p1a}  \\[1.5mm]
			&  \|\mathcal{S}_{\theta} u-u\|_{H_*^{k}(\Omega_{T})}\lesssim_m \theta^{k-j}\|u\|_{H_*^{j}(\Omega_{T})}
			&& \textrm{for  \;}0\leq k\leq j \leq m,   \label{smooth.p1b} \\
			&  \left\|\frac{\d}{\d \theta}\mathcal{S}_{\theta} u\right\|_{H_*^{k}(\Omega_{T})}
			\lesssim_m \theta^{k-j-1}\|u\|_{H_*^{j}(\Omega_{T})}
			&\qquad&\textrm{for \;}  k,j\geq 0,    \label{smooth.p1c}
		\end{alignat}
	\end{subequations}
	where $(k-j)_+:=\max\{0,\, k-j \}$.
	Moreover, there are two families of smoothing operators (still denoted by $\mathcal{S}_{\theta}$) acting respectively on $\mathscr{F}^6(\Omega_{T})$ and $\mathscr{F}^6(\Sigma_T)$ that satisfy  the properties  \eqref{smooth.p1} with corresponding norms $\|\cdot\|_{H^{s}(\Omega_{T})}$ and   $\|\cdot\|_{H^{s}(\Sigma_T)}$ with $s\geq 0$ (not necessarily an integer).
\end{proposition}

\begin{align} \label{source}
	\sum_{k=0}^{{n}} f_k^+ +\mathcal{S}_{\theta_{{n}}}E_{{n}}^+=\mathcal{S}_{\theta_{{n}}}f^a,\quad
	\sum_{k=0}^{{n}}f_k^{-}+\mathcal{S}_{\theta_{{n}}}{E}_{{n}}^{-}=0,\quad
	\sum_{k=0}^{{n}}g_k +\mathcal{S}_{\theta_{{n}}}\widetilde{E}_{{n}} =0,
\end{align}

\begin{lemma} \label{lem:sum2}
	Let $\widetilde{\alpha}={\alpha }+4$ and ${\alpha }\geq 13$.
	If $\theta_0\geq 1$ is sufficiently  large and $\varepsilon,T>0$ are small enough, then	
	\begin{align} \label{es.sum2}
		\|E_{N}^+\|_{H_*^{{\alpha }+2}(\Omega_{T})}+\|E_{N}^-\|_{H^{{\alpha }+3}(\Omega_{T})}
		\lesssim \varepsilon^2 \theta_{{N}},\quad
		\| \widetilde{E}_{{N}} \|_{H^{{\alpha }+5/2}(\Sigma_{T})}\lesssim \varepsilon^2.
	\end{align}
\end{lemma}

	\begin{align} \label{es.sum1a}
		\left\{
		\begin{aligned}
			&	\|e_{n}^+ \|_{H_*^s(\Omega_{T})}+\|e_{n}^- \|_{H^{s+1}(\Omega_{T})}
			\lesssim \varepsilon^2 \theta_n^{\varsigma_4(s)-1}\varDelta_n,\\[1mm]
			&	\|\tilde{e}_n\|_{H^{s+1/2}(\Sigma_T)}
			\lesssim \varepsilon^2 \theta_n^{\varsigma_2(s)-1}\varDelta_n,
		\end{aligned}
		\right.
	\end{align}	

\begin{lemma} \label{lem:source}
	Let $\widetilde{\alpha}={\alpha }+4$ and ${\alpha }\geq 13$.
	If $\theta_0\geq 1$ is large enough and $\varepsilon,T>0$ are sufficiently small, then
	\begin{align}
		&
		\|f_{{N}}^+\|_{H_*^s(\Omega_{T})}+\|f_{{N}}^-\|_{H^{s+1}(\Omega_{T})}
		\lesssim \varDelta_{{N}}\left(\theta_{{N}}^{s-{\alpha }-2} (\|f^a\|_{H_*^{\alpha+1}(\Omega_{T})}
		+\varepsilon^2 )  +\varepsilon^2\theta_{{N}}^{\varsigma_4(s)-1}\right),
		\nonumber \\[0.5mm] %\label{es.fl} \\
		&
		\| g_{{N}} \|_{H^{s+3/2}(\Sigma_{T}) }
		\lesssim  \varepsilon^2 \varDelta_{{N}}\left(\theta_{{N}}^{s-{\alpha }-2}
		+\theta_{{N}}^{\varsigma_2(s+1)-1}\right)\nonumber
		\qquad\qquad \textrm{for }s=6,7,\ldots,\widetilde{\alpha}+1.
	\end{align}
\end{lemma}

\begin{proof}
	It follows from Proposition \ref{pro:smooth} and estimates \eqref{es.sum1a} that
	\begin{align*}
		&	\big\|\mathcal{S}_{\theta_{{N}}} e_{N-1}^+ \big\|_{H_*^s(\Omega_{T})}
		\lesssim 	\big\|  e_{N-1}^+ \big\|_{H_*^s(\Omega_{T})}
		\lesssim \varepsilon^2 \theta_N^{\varsigma_4(s)-1}\varDelta_N
		\quad&& \textrm{if } 6\leq s\leq \widetilde{\alpha}-2,\\
		&	\big\|\mathcal{S}_{\theta_{{N}}} e_{N-1}^+ \big\|_{H_*^s(\Omega_{T})}
		\lesssim \theta_{{N}}^{s-\widetilde{\alpha}+2}	
		\big\|  e_{N-1}^+ \big\|_{H_*^{\widetilde{\alpha}-2}(\Omega_{T})}
		\lesssim \varepsilon^2 \theta_N^{\varsigma_4(s)-1}\varDelta_N
		\quad &&\textrm{if }  s\geq \widetilde{\alpha}-2.
		%	+\|e_{n}^- \|_{H^{s+1}(\Omega_{T})}\lesssim \varepsilon^2 \theta_n^{\varsigma_4(s)-1}\varDelta_n,\quad 	\|\tilde{e}_n\|_{H^{s}(\Sigma_T)}	\lesssim \varepsilon^2 \theta_n^{\varsigma_2(s)-1}\varDelta_n,
	\end{align*}
	Similarly, for all $s\geq 6$, we have
	\begin{align*}
		\big\|\mathcal{S}_{\theta_{{N}}} e_{N-1}^- \big\|_{H_*^{s+1}(\Omega_{T})}
		\lesssim \varepsilon^2 \theta_N^{\varsigma_4(s)-1}\varDelta_N,
		\quad 	
		\big\|\mathcal{S}_{\theta_{{N}}} \tilde{e}_{N-1}  \big\|_{H_*^{s+3/2}(\Sigma_{T})}
		\lesssim \varepsilon^2 \theta_N^{\varsigma_2(s+1)-1}\varDelta_N.
	\end{align*}
	By virtue of  identities \eqref{source}, Proposition \ref{pro:smooth}, and Lemma \ref{lem:sum2}, we infer
	\begin{align*}
		\big\|f_N^+\big\|_{H_*^{s}(\Omega_{T})}
		&\lesssim
		\big\|(\mathcal{S}_{\theta_{{N}}}-\mathcal{S}_{\theta_{{N}-1}})f^a
		-(\mathcal{S}_{\theta_{{N}}}-\mathcal{S}_{\theta_{{N}-1}})E_{{N}-1}^+
		-\mathcal{S}_{\theta_{{N}}}e_{{N}-1}^+\big\|_{H_*^{s}(\Omega_{T})} \\
		&
		\lesssim \varDelta_{{N}}
		\theta_{{N}}^{s-{\alpha }-2}\big(\big\|f^a\big\|_{H_*^{\alpha +1}(\Omega_{T})}
		+\theta_{{N}}^{-1}\big\|E_{{N}-1}^+\big\|_{H_*^{\alpha +2}(\Omega_{T})} \big)
		+\varepsilon^2 \theta_N^{\varsigma_4(s)-1}\varDelta_N,\\[1mm]
		\big\|f_{{N}}^-\big\|_{H^{s+1}(\Omega_{T})}
		&\lesssim \big\|(\mathcal{S}_{\theta_{{N}}}-\mathcal{S}_{\theta_{{N}-1}})E_{{N}-1}^-\big\|_{H^{s+1}(\Omega_{T})}
		+\big\|\mathcal{S}_{\theta_{{N}}}e_{{N}-1}^-\big\|_{H^{s+1}(\Omega_{T})}\\
		&\lesssim
		\varDelta_{{N}}
		\theta_{{N}}^{s-{\alpha }-3}\big\|E_{{N}-1}^-\big\|_{H^{\alpha+3}(\Omega_{T})}
		+\varepsilon^2 \theta_N^{\varsigma_4(s)-1}\varDelta_N,\\[1mm]
		\big\|g_{{N}} \big\|_{H^{s+3/2}(\Sigma_{T})}
		&\lesssim \big\|(\mathcal{S}_{\theta_{{N}}}-\mathcal{S}_{\theta_{{N}-1}})\widetilde{E}_{{N}-1}\big\|_{H^{s+3/2}(\Sigma_{T})}
		+\big\|\mathcal{S}_{\theta_{{N}}}\tilde{e}_{{N}-1}\big\|_{H^{s+3/2}(\Sigma_{T})}\\
		&\lesssim
		\varDelta_{{N}}
		\theta_{{N}}^{s-{\alpha }-2}\big\|\widetilde{E}_{{N}-1}\big\|_{H^{\alpha+5/2}(\Sigma_{T})}
		+\varepsilon^2 \theta_N^{\varsigma_2(s+1)-1}\varDelta_N,
	\end{align*}
	which together with Lemma \ref{lem:sum2} completes the proof of this lemma.
\end{proof}

\end{document}
